\section{Exponential tails for the genuine return record}
\label{sec:exponential-returns}
We now use the spectral theory of the collision map, in addition to the
geometric and ergodic inputs used above. The operator in this section is
not the operator of the induced map. A smooth killing weight on the
collision space gives a scalar estimate for avoiding the actual return
section. In particular no multiplication by its discontinuous indicator
on an anisotropic space will be needed.

\subsection{The collision-space input}
We use the strong and weak distribution spaces of Demers--Zhang
\cite[Section 3.3]{DZ}. Their properties needed here are the local uniform
spectral gap and power bounds for unforced dispersing collision maps,
including small changes of the scatterers and their boundary lengths
\cite[Theorems 2.1--2.6 and Remark 2.7(b)]{DZ}. The spectral gap for an
individual unforced Lorentz map is the input from \cite{DZ11} used there.
Here is the precise consequence, together with the checks for this family.

\begin{lemma}[Local uniform collision spaces]
\label{lem:collision-spectral-input}
There is a finite cover of $I$ by parameter neighborhoods $I_j$. On each
$I_j$ there is a Banach space $\mathcal B_j$ of distributions on a common
collision space, a continuous mass functional $\ell_j$, and transfer
operators $\mathcal L_R$, such that
\begin{equation}\label{eq:collision-split}
 \mathcal L_R=\Pi_R+\mathcal N_R,\quad
 \Pi_R h=\ell_j(h)\nu,\quad
 \Pi_R\mathcal N_R=\mathcal N_R\Pi_R=0,\quad
 \|\mathcal N_R^k\|_{\mathcal B_j}\le C\vartheta^k
 \quad(k\ge0),
\end{equation}
where $0<\vartheta<1$ and $C<\infty$ can be chosen for all the finitely
many neighborhoods. The norms of $\nu$, $\ell_j$ and $\mathcal L_R$ are
uniformly bounded. Multiplication by a smooth function $g$ satisfies
\begin{equation}\label{eq:smooth-multiplier}
 \|g h\|_{\mathcal B_j}\le C\|g\|_{C^2}\|h\|_{\mathcal B_j}.
\end{equation}
The transfer convention is $(\mathcal L_R h)(f)=h(f\circ T_R)$.
No assertion about the unbounded induced twists is included in this lemma.
\end{lemma}
\begin{proof}
For the collision-space construction use coordinates $(\alpha,\varphi)$,
where $p=\sin\varphi$. The invariant probability is the same for all
parameters: $\dd\nu=(4\pi)^{-1}\cos\varphi\dd\alpha\dd\varphi$.
We do not conjugate a distribution space by the inverse of
$p=\sin\varphi$ at grazing. At a fixed $R_0$, use the boundary coordinate
$r_0=R_0\alpha$. The parametrization of the nearby circle is then
$r_0\mapsto R n_{r_0/R_0}$. It is a $C^3$-bounded deformation, $C^2$-close
to the original parametrization, and its constant speed $R/R_0$ tends to
one. This is precisely the boundary-length reparametrization permitted in
\cite[Remark 2.7(b)]{DZ}.

The physical gap between different disks is at least $3/50$. The
curvatures belong to $[100/47,20/9]$, and the parametrized boundaries have
uniform $C^3$ bounds. The finite horizon is uniform on $I$: a line with
no disk intersection at the smallest radius would be an infinite free
line in that table, which has no corridor; compactness of position and
direction then bounds the free length, and increasing the radius cannot
increase it. If a rectangular fundamental domain is desired, use the
two-scatterer cover with periods $(1,0)$ and $(0,\sqrt3)$ and disk centers
$0,b$. Its deck translation commutes with the collision map. Restricting
to the deck-invariant distributions gives the original table. Thus the
local billiard estimates are applied without an affine distortion of
specular reflection.

The uniform geometric hypotheses of \cite[Theorem 2.5]{DZ} and the
perturbation hypothesis of its Theorem 2.6 are satisfied. The unforced
collision map at each $R_0$ has the simple eigenvalue one and a spectral
gap. The strong/weak perturbation theorem, specifically the complementary
power bound in \cite[Theorem 2.1(3)]{DZ}, gives \eqref{eq:collision-split}
on a neighborhood of $R_0$. The eigenprojection has the displayed form
because the invariant probability is $\nu$ and mass is preserved. A
finite cover of the compact parameter interval makes the constants
uniform. This uses local common spaces and does not assert that all
arbitrary inducing schemes share a Banach space.

For completeness, smooth multiplication follows directly from the
stable-test and unstable-comparison norms defining these spaces. On a
single admissible stable curve, multiplying its test function by $g$
changes its H\"older norm by at most $C\|g\|_{C^1}$. For two matched
curves at distance $\epsilon$, split the difference of the weighted tests
into the old test difference and the difference of the two restrictions
of $g$. After graph matching, the latter has $C^q$ norm at most
$C\|g\|_{C^2}\epsilon^{1-q}$ by interpolation between its $C^0$ and
$C^1$ bounds. The defining exponents satisfy $\beta<1-q$, so division
by $\epsilon^\beta$ is bounded for $\epsilon\le1$. The stable part of the
norm controls this second term and the unstable part controls the first.
The weak norm has the same single-curve multiplication bound. Density of
the defining smooth distributions extends these bounds to
$\mathcal B_j$, proving \eqref{eq:smooth-multiplier}. The coordinate
changes above and the finite cover have uniformly bounded derivatives.
\end{proof}
\subsection{A smooth avoidance estimate}
Choose once and for all $\chi\in C^\infty(M)$ such that
\begin{equation}\label{eq:fixed-killing-weight}
 0\le\chi\le1,\quad
 \supp\chi\Subset\{ |\alpha|<1/400,\ |p|<1/400\},\quad
 a_\chi:=\int\chi\dd\nu>0.
\end{equation}
The indicated rectangle is inside the square of $Y_R$ centered at
$(0,0)$ for every $R$; hence it is inside $Y_R^*$. Its support is away
from grazing, so it is smooth also in the collision coordinates used in
Lemma~\ref{lem:collision-spectral-input}. This is a proof weight, not an
extra field in the collision record.

\begin{proposition}[Uniform soft-killing estimate]
\label{prop:soft-killing}
There exist $z_0,c,C>0$ such that for every $R\in I$ and $k\ge0$,
\begin{equation}\label{eq:soft-killing}
 \int_M\exp\left(-z_0\sum_{j=0}^{k-1}\chi\circ T_R^j\right)\dd\nu
       \le Ce^{-ck}.
\end{equation}
\end{proposition}
\begin{proof}
On a local collision space set
$\mathcal L_{R,z}=\mathcal L_R M_{e^{-z\chi}}$ for complex $z$.
The multiplier bound proves analyticity in the strong operator norm and,
uniformly in $R$, $\|\mathcal L_{R,z}-\mathcal L_R\|\le C|z|$ near
zero. Equation~\eqref{eq:collision-split} bounds the unperturbed
resolvents uniformly on a circle around one and on a second circle
$|w|=\vartheta_1$ with $\vartheta<\vartheta_1<1$, chosen disjointly.
A uniformly convergent resolvent Neumann series therefore defines a
rank-one analytic projection and a single eigenvalue $\lambda_R(z)$
near one for all $|z|<z_*$. The complementary powers are bounded by
$C\vartheta_1^k$, after decreasing $z_*$ if necessary.

Since $\ell_j(\nu)=1$, differentiation at zero gives
\[
 \lambda_R(0)=1,\qquad
 \lambda_R'(0)=-\ell_j(\mathcal L_R(\chi\nu))=-a_\chi.
\]
The fixed complex disk and the uniform resolvent bounds give a uniform
bound on the second derivative of $\lambda_R$ on a smaller disk, by
Cauchy's formula. Thus
$|\lambda_R(z)-1+a_\chi z|\le C_2|z|^2$ uniformly in $R$.
For real small $z$ the eigenvalue is real, by conjugation and uniqueness,
and positive, by its closeness to one. Choose a fixed $z_0>0$ such that
$C_2z_0\le a_\chi/2$ and $z_0<z_*$. Then
$0<\lambda_R(z_0)\le1-a_\chi z_0/2$. The spectral decomposition gives
$\|\mathcal L_{R,z_0}^k\|\le Ce^{-ck}$. Only finitely many collision
spaces occur, so the same $z_0,c,C$ work on all of $I$.

Iteration of the transfer identity, first for densities and then as
bounded functionals, yields
\[
 \ell_j(\mathcal L_{R,z_0}^k\nu)
 =\int\prod_{j=0}^{k-1}e^{-z_0\chi(T_R^j x)}\dd\nu(x).
\]
The norms of $\ell_j$ and $\nu$ are bounded. This proves the asserted
scalar estimate. In particular the argument does not apply a local
large-deviation formula outside its stated neighborhood of the mean.
\end{proof}

\begin{theorem}[Exponential moments of the actual return block]
\label{thm:exponential-return}
Let
$Q_R=r_R^*+\tau_R^*+|\kappa_R^*|$ on the genuine first-return section.
There are constants $c_1,a_1,C_1>0$, independent of $R\in I$, such that
\begin{equation}\label{eq:exponential-return}
 \nu_R^*(r_R^*>k)\le C_1e^{-c_1k}\quad(k\ge0),\qquad
 \int e^{a_1 Q_R}\dd\nu_R^*\le C_1.
\end{equation}
For $Q_{n,R}=\sum_{j<n}Q_R\circ(F_R^*)^j$ and every $n\ge1$,
\begin{equation}\label{eq:block-exponential-moment}
 \int\exp(a_1 Q_{n,R}/n)\dd\nu_R^*\le C_1,\qquad
 \nu_R^*(N_{n,R}>L)\le C_1e^{-a_1L/n}.
\end{equation}
These are estimates for the original deterministic mechanical records.
\end{theorem}
\begin{proof}
If $x\in Y_R^*$ and $r_R^*(x)>k$, none of $T_Rx,\ldots,T_R^kx$
belongs to $Y_R^*$. Every term of $\sum_{j=1}^k\chi\circ T_R^j$ is
therefore zero. Positivity and invariance of $\nu$ imply
\[
 \nu_R^*(r_R^*>k)
 \le c_*^{-1}\int_M e^{-z_0\sum_{j=1}^k\chi\circ T_R^j}\dd\nu
 \le c_*^{-1}Ce^{-ck}.
\]
This comparison uses neither the regularity of $\one_{Y_R^*}$ as a
multiplier nor an independent-return construction.

Uniform finite horizon and the uniform bound on a single-flight lattice
step give $Q_R\le C_0 r_R^*$ for a fixed $C_0$. Summation of the
exponential tail of the positive integer $r_R^*$ gives a uniform
exponential moment of $Q_R$ for $a_1>0$ sufficiently small. Enlarge
$C_1$ to include the first bound. Finally H\"older's inequality with
$n$ equal exponents and invariance of $F_R^*$ give
\[
 \int\prod_{j<n}\exp\bigl(a_1Q_R\circ(F_R^*)^j/n\bigr)\dd\nu_R^*
 \le\prod_{j<n}\left(\int e^{a_1Q_R\circ(F_R^*)^j}\dd\nu_R^*\right)^{1/n}
 \le C_1.
\]
Since $N_{n,R}\le Q_{n,R}$, exponential Markov inequality proves the
last assertion. No independence is used in this step either.
\end{proof}

\begin{corollary}[All finite moments under parameter variation]
\label{cor:all-moment-continuity}
For fixed $n$ and $1\le p<\infty$, the map
$R\mapsto\widetilde J_{n,R}$ is continuous into $L^p(\nu;\R^4)$.
Every polynomial moment of a fixed finite list of return records is
finite and continuous in $R$. In particular every finite-lag covariance
of the induced record is continuous.
\end{corollary}
\begin{proof}
Almost-everywhere convergence is Lemma~\ref{lem:finite-return-stability}.
Equation~\eqref{eq:block-exponential-moment}, the fixed section mass and
$|J_{n,R}|\le Q_{n,R}$ make all fixed powers uniformly integrable.
Truncation followed by dominated convergence proves $L^p$ convergence
along every convergent parameter sequence. A finite list of records can
be expressed as differences of cumulative records. H\"older's inequality
and the same exponential moment bound give uniform integrability of each
polynomial in that list. Truncation proves its continuity. This proves
finite-lag, not yet infinite Green--Kubo, covariance continuity.
\end{proof}
